%2multibyte Version: 5.50.0.2953 CodePage: 1251
\documentclass[12pt]{article}%
\usepackage{amsmath,amsthm,enumerate,graphicx,microtype,amssymb,setspace,verbatim,multirow}
\usepackage[margin= 1in]{geometry}
\usepackage[font={footnotesize}]{caption,subfig}
\usepackage{caption}
\usepackage{setspace}
\usepackage{natbib}
\usepackage{subfig}
\usepackage{amsmath}
\usepackage{amsfonts}
\usepackage{amssymb}
\usepackage{graphicx}
\usepackage{epstopdf}
\usepackage{dsfont}
\usepackage{bbm}
\usepackage{setspace}
\usepackage{soul}
\usepackage{color}
\usepackage{chngcntr}
\usepackage{todonotes}%
\setcounter{MaxMatrixCols}{30}
%TCIDATA{OutputFilter=latex2.dll}
%TCIDATA{Version=5.50.0.2953}
%TCIDATA{Codepage=1251}
%TCIDATA{LastRevised=Wednesday, November 08, 2023 18:51:01}
%TCIDATA{<META NAME="GraphicsSave" CONTENT="32">}
%TCIDATA{<META NAME="SaveForMode" CONTENT="1">}
%TCIDATA{BibliographyScheme=Manual}
%BeginMSIPreambleData
\providecommand{\U}[1]{\protect\rule{.1in}{.1in}}
%EndMSIPreambleData
\providecommand{\U}[1]{\protect\rule{.1in}{.1in}}
\onehalfspacing
\renewcommand{\labelenumiii}{(\roman{enumiii})}
\DeclareMathOperator{\Cov}{Cov}
\DeclareMathOperator{\Var}{Var}
\DeclareMathOperator{\diag}{diag}
\setlength{\parindent}{1 cm}
\newtheorem*{Prop1}{Proposition 1}
\newtheorem*{Prop2}{Proposition 2}
\newtheorem*{Prop3}{Proposition 3}
\newtheorem{Prop}{Proposition}
\newtheorem{definition}{Definition}
\newtheorem{lemma}{Lemma}
\newtheorem{corollary}{Corollary}
\newtheorem{corollary2}{Corollary 2}
\newcommand{\chg}[2]{{\st{#1}} {\color{red}{#2}}}
\newcommand{\distas}[1]{\mathbin{\overset{#1}{\kern\z@\sim}}}
\newsavebox{\mybox}\newsavebox{\mysim}
\newcommand{\distras}[1]{	\savebox{\mybox}{\hbox{\kern3pt$\scriptstyle#1$\kern3pt}}	\savebox{\mysim}{\hbox{$\sim$}}	\mathbin{\overset{#1}{\kern\z@\resizebox{\wd\mybox}{\ht\mysim}{$\sim$}}}}
\makeatother
\renewcommand{\thefigure}{\Roman{figure}}
\renewcommand{\thetable}{\Roman{table}}
\pagenumbering{arabic} \onehalfspacing
\begin{document}

\title{An economic model of reinforcement learning}
\maketitle

.\bigskip

\section{Definitions}

Policy is a mapping from states to probabilities over action
\[
\pi(a|s)
\]
Q-function: State-action values of following a policy $\pi(a|s)$ in the
future, and implement a specific $a_{t}$ today%

\[
Q_{\pi}(s_{t},a_{t})=r(s_{t},a_{t})+\beta E_{t}Q_{\pi}(s_{t+1},\pi(s_{t+1}))
\]


Value function: value in state $s$ of following policy $\pi(a|s)$
\begin{equation}
V_{\pi}(s)=\sum_{a}\pi(a|s)Q_{\pi}(s,a) \label{v_pi_s}%
\end{equation}


\section{Bayesian resource-rational framework of learning}

\subsection{Object the agent is learning about}

We assume the agent is trying to learn the true, full-information optimal
Q-function $Q_{\pi^{*}}(s,a)$ where the optimal policy $\pi^{*}(a|s)$
satisfies the Bellman equation%

\[
V^{*}(s_{t}) = \max_{a_{t}} r(s_{t},a_{t}) + \beta E_{t}( V(s_{t+1}))
\]


The optimal policy function $\pi^{*}(a|s)$ defines the optimal Q-function:
\[
Q_{\pi^{*}}(s_{t},a_{t})=r(s_{t},a_{t})+\beta E_{t}Q_{\pi^{*}} (s_{t+1}%
,\pi^{*}(s_{t+1}))
\]


The agent will learn about the optimal Q-value function $Q_{\pi^{*}}(s,a)$
defined above.

\subsection{Policies}

Each period the agent enters with some beliefs about the Q-function, those
beliefs follow a Gaussian Process distribution so that conditional on the
$t-1$ information set $\mathcal{I}_{t-1}$ (which is what the agent enters time
$t$ with)%

\[
Q_{\pi^{*}}(s,a) | \mathcal{I}_{t-1} \sim GP( \widehat Q_{t-1}(s,a),
\widehat\Sigma_{t-1}( s,a,s^{\prime},a^{\prime}))
\]


In period $t$, the agent follows the following policy function that defines
his actions. Conditional on the (updated) information set $\mathcal{I}_{t}$,
the probability of selecting action $a$ given state $s_{t}$ is
\[
\pi_{t}(s_{t})=\pi(\widehat{Q}_{t}(s_{t},a),\theta_{t}(s_{t}))=\frac
{\exp(\theta_{t}(s_{t})\widehat{Q}_{t}(s_{t},a))}{\sum_{a^{\prime}}\exp
(\theta_{t}(s_{t})\widehat{Q}_{t}(s_{t},a^{\prime}))}%
\]


Here $\widehat{Q}_{t}(s_{t},a)$ is the conditional expectation of the Q-value
function and $\theta_{t}(s_{t})$ is a temperature parameter. We can treat the
latter as a constant%
\[
\theta_{t}(s)=\theta,\text{ for all (}s,t)
\]
or, more interestingly, have it take the following form:
\[
\theta_{t}(s)=\theta_{0}\frac{\Sigma_{V,0}(s)}{\Sigma_{V,t-1}(s)}%
\]
where $\Sigma_{V}(s)$ will be a measure of perceived uncertainty, which we
will define below.

Also define the 'greedy' policy as
\begin{align*}
\pi_{t-1}^{greedy}(a_{t}  &  =\widehat{a}_{t-1}^{\max}(s_{t})|s_{t})=1\\
\pi_{t-1}^{greedy}(a_{t}  &  \neq\widehat{a}_{t-1}^{\max}(s_{t})|s_{t})=0\\
\widehat{a}_{t-1}^{\max}(s_{t})  &  =\arg\max_{a}\widehat{Q}_{t-1}(s_{t}%
,a_{t})
\end{align*}


\subsection{Updating}

For any state $s,$ compute the expected return%
\[
\widehat{V}_{t-1}(s)=\sum_{a}\pi_{t-1}(a_{t}=a|s)\widehat{Q}_{t-1}(s,a_{t})
\]


This is a variable measurable in the time $t-1$ information set, hence it is
known at the beginning of time $t$. It is the expected value of the following
random variable:%

\[
V_{t-1}^{true}(s)\equiv\sum_{a}\pi_{t-1}(a_{t}=a|s)Q_{\pi}^{\ast}(s,a_{t})
\]
The policy $\pi_{t-1}(a|s)$ is known as of time $t$, and the randomness comes
from the uncertainty over the unknown optimal $Q_{\pi^{*}}(s,a).$

Based on prior information, this random variable has a distribution%
\begin{align*}
V_{t-1}^{true}(s_{t})  &  \sim N(\widehat{V}_{t-1}(s_{t}),\Sigma_{V,t-1}%
(s_{t}))\\
\Sigma_{V,t-1}(s_{t})  &  = Var_{t-1}( \sum_{a}\pi_{t-1}(a|s_{t})Q_{\pi^{*}%
}(s_{t},a)) = \pi_{t-1}(a|s_{t})^{\prime}\widehat{\Sigma}_{t-1}(a,s_{t}%
,a^{\prime},s_{t}) \pi_{t-1}( a^{\prime}|s_{t})
\end{align*}
where $\widehat{\Sigma}_{t-1}(a,s,a^{\prime},s)$ is a two-dimensional matrix,
which evaluates the variance-covariance of $Q_{\pi^{*}}(s,a)$ at a fixed
$s_{t}$. So then $\pi_{t-1}(a|s_{t})^{\prime}\widehat{\Sigma}_{t-1}%
(a,s_{t},a^{\prime},s_{t}) \pi_{t-1}( a^{\prime}|s_{t})$ is the typical
sandwich expression for the variance of the linear combination $\sum_{a}%
\pi_{t-1}(a|s_{t})\widehat{Q}_{t-1}(s_{t},a)$.

%we define $\sigma_{Q,t-1}^{2}(s,a) = \widehat{\Sigma}_{t-1}(s,a,s,a) = Var_{t-1}( Q_{\pi^*}(s,a))$ as the posterior variance over $Q$ at the particular input of $(s,a)$. It is computed by evaluating the variance-covariance function at $(s,a,s,a)$ (i.e. it is on the diagonal).


\subsubsection{Deliberation (i.e. model-based reasoning)}

\paragraph{Resource-rationality, conditional trigger of reasoning}

The agent looks at $\Sigma_{V,t-1}(s_{t})$ and decides whether to engage in
additional reasoning to reduce this uncertainty. The agent's objective
function in deciding whether to engage in extra reasoning is a tracking
problem of the form
\[
\min_{\Sigma_{V,t}(s_{t})}E\left(  \widehat{V}_{t}(s_{t})-V_{t-1}^{true}%
(s_{t})\right)  ^{2}+\kappa\ln\left(  \frac{\Sigma_{V,t-1}(s_{t})}%
{\Sigma_{V,t}(s_{t})}\right)
\]
where $\kappa$ is the marginal cost of reducing further uncertainty through
reasoning. This objective is then%
\[
\min_{\widetilde{\Sigma}_{V,t}(s_{t})} \widetilde{\Sigma}_{V,t}(s_{t}%
)+\kappa\ln\left(  \frac{\Sigma_{V,t-1}(s_{t})}{\widetilde{\Sigma}_{V,t}%
(s_{t})}\right)  .
\]
where
\[
\widetilde{\Sigma}_{V,t}(s_{t})= Var(\sum_{a}\pi_{t-1}(a|s_{t}) Q_{\pi^{*}}
(s,a) | \mathcal{I}_{t-1}, \eta_{t}^{R})
\]
where $Var( Q_{\pi^{*}} (s,a) | \mathcal{I}_{t-1}, \eta_{t}^{R})$ is the
updated conditional uncertainty after observing the new signal $\eta_{t}^{R}$,
which we define below. The optimal information choice is%

\begin{align}
\label{eq: opt_info}\widetilde{\Sigma}_{V,t} = \min\{ \kappa, \Sigma
_{V,t-1}(s_{t}) \}
\end{align}


\subsubsection{Mental simulation, model-based learning}

Each period $t$, internal deliberation produces a noisy signal $\eta_{t}%
^{R}(s_{t},a_{t})$%

\[
\eta_{t}^{R}=\sum_{a}\pi_{t-1}(a|s_{t})Q_{\pi^{*}}(s_{t},a )+\varepsilon
_{t}^{R}%
\]
where $\varepsilon_{t}^{R} \sim N(0, \sigma_{\eta^{R},t}^{2})$, where the
variance of the noise follows from the above optimal information choice in
(\ref{eq: opt_info}):%

\[
\sigma_{\eta^{R},t}^{2}=
\begin{cases}
\frac{\kappa{\Sigma}_{V,t-1}(s_{t})}{ {\Sigma}_{V,t-1}(s_{t})-\kappa} & \text{
, if } {\Sigma}_{V,t-1}(s_{t}) > \kappa.\\
\infty & \text{ , if } {\Sigma}_{V,t-1}(s_{t}) \leq\kappa.
\end{cases}
\]


Then use that signal to update beliefs%

\[
\widehat{Q}_{t}^{R}(s,a)=\widehat{Q}_{t-1}(s,a)+\alpha_{t}^{R}(s,a)\left(
\eta_{t}^{R}-\sum_{a^{\prime}}\pi_{t-1}(a^{\prime}|s_{t})\widehat{Q}%
_{t-1}(s_{t},a^{\prime})\right)
\]


where we can derive $\alpha_{t}^{R}(s,a)$ as follows:%
\begin{align*}
\alpha_{t}^{R}(s,a)  &  =\frac{Cov_{t-1}(Q_{\pi^{*}}(s,a),\sum_{a^{\prime}}%
\pi_{t-1}(a^{\prime}|s_{t})Q_{\pi^{*}}(s_{t},a^{\prime}))}{Var_{t-1}(
\sum_{a^{\prime}}\pi_{t-1}(a^{\prime}|s_{t})Q_{\pi^{*}}(s_{t},a^{\prime
})+\varepsilon_{t}^{R})}\\
&  =\sum_{a^{\prime}}\pi_{t-1}(a^{\prime}|s_{t})\frac{Cov(Q_{\pi^{*}%
}(s,a),Q_{\pi^{*}}(s_{t},a^{\prime}))}{\Sigma_{V,t-1}(s_{t}) + \sigma
_{\eta^{R},t}^{2}}%
\end{align*}


Similarly, posterior variances and etc. of $Q_{\pi}^{\ast}(s,a)$ update as
well at all $(s,a)$ pairs.

Cost of reasoning is the reduction of uncertainty (at optimal reasoning)
\[
c_{t-1/t}(s_{t})\equiv\kappa\ln\left(  \frac{\Sigma_{V,t-1}(s_{t})}%
{\Sigma_{V,t} (s_{t})}\right)
\]
By optimal reasoning this is only a function of previous uncertainty,
evaluated at current state%

\[
c_{t-1}(s_{t})=
\begin{cases}
\kappa\ln(\frac{\Sigma_{V,t-1}(s_{t})}{\kappa}) & \text{ , if } {\Sigma
}_{V,t-1}(s_{t}) > \kappa.\\
0 & \text{ , if } {\Sigma}_{V,t-1}(s_{t}) \leq\kappa.
\end{cases}
\]


Therefore experienced reward at any time $t$ is
\[
r(s_{t},a_{t})\equiv u(s_{t},a_{t})-c(s_{t})
\]
where $u(s_{t},a_{t})$ might be physical utility.

\subsubsection{Model-free learning}

Each period the agent also observes the utility he just experienced at the end
of last period, $u(s_{t-1},a_{t-1})$. He can use that observed flow utility
value to update beliefs about $Q_{\pi^{*}}(s,a)$ as well. We think of this
learning ``by experience'' as free, so it will not be subject to a cognitive
cost. Still, it is limited as we can see below.

It is most intuitive to write this ``experience'' signal as
\[
\eta_{t}^{Q} = r(s_{t-1}, a_{t-1}) + \beta\sum_{a} \sum_{s_{t}} \pi
_{t-1}^{greedy} (a |s_{t}) \widehat Q_{t-1}(s_{t},a) F(s_{t} | s_{t-1})
\]
where $F(s_{t} | s_{t-1})$ is the law of motion for the state (that could also
depend on the choice $a_{t-1}$, but for simplicity we suppress this notation).
So the signal $\eta_{t}^{Q}$ is sort of like the expected value of $Q(s_{t-1},
a_{t-1})$ from the perspective of time $t-1$. Note: we use the greedy policy
because we're learning about the (unknown) Q-value derived under the (unknown)
greedy policy.

The new ``information'' is in the fact that we will now impose the Q-value equation%

\[
Q_{\pi^{*}}(s_{t},a_{t})=r(s_{t},a_{t})+\beta E_{t}Q_{\pi^{*}} (s_{t+1}%
,\pi^{*}(s_{t+1}))
\]
which up to this point the agent has never used directly (only abstractly, in
reduced form in the noisy signals $\eta_{\eta^{R},t}$. Important: we do not
assume this equation is in the agent's information set. Assuming knowledge of
this equation automatically means we are assuming the agent knows $Q_{\pi^{*}%
}(s,a)$, since the optimal value function is defined by that equation and is
just a direct one-to-one map from it.

Note: we can just rewrite that as something that makes explicit the greedy
policy
\[
Q_{\pi^{*}}(s_{t},a_{t})=r(s_{t},a_{t})+\beta E_{t}\max_{a^{\prime}}Q_{\pi
^{*}}(s_{t+1},a^{\prime})
\]


So we are not going to give the agent full knowledge of this equation, we'll
just let it use once in a way. Specifically, we will use it to express the
experience signal as%

\begin{align*}
\eta_{t}^{Q} = u(s_{t-1}, a_{t-1})  &  = r(s_{t-1}, a_{t-1}) + \beta\sum_{a}
\sum_{s_{t}} \pi_{t-1}^{greedy} (a |s_{t}) \widehat Q_{t-1}(s_{t},a) F(s_{t} |
s_{t-1})\\
&  = Q^{*}(s_{t-1}, a_{t-1}) + \beta\sum_{a} \sum_{s_{t}} \pi_{t-1}^{greedy}
(a |s_{t}) F(s_{t} | s_{t-1}) \left(  \widehat Q_{t-1}(s_{t},a) - Q_{\pi^{*}%
}(s_{t}, a) \right)
\end{align*}


The agent uses that to update beliefs {\small
\[
\hspace{-1cm} \widehat{Q}_{t}(s,a)=\widehat{Q}_{t-1}(s,a)+\alpha_{t}%
^{Q}(s,a)\left[  \eta_{t}^{Q}- \widehat{Q}_{t-1} (s_{t-1},a_{t-1}) \right]
\]
}

where%

\[
\alpha_{t}^{Q}(s,a) = \frac{ Cov_{t-1}( \eta_{t}^{Q}, Q^{*}(s,a))}{Var_{t-1}(
\eta_{t}^{Q})}
\]
with
\begin{align*}
\hspace{-1cm} Cov_{t-1}( \eta_{t}^{Q}, Q^{*}(s,a))  &  = Cov_{t-1}(
Q^{*}(s_{t-1}, a_{t-1}) - \beta\sum_{a} \sum_{s} \pi_{t-1}^{greedy}(a |s_{t})
F(s_{t}| s_{t-1}) Q^{*}(s_{t},a_{t}), Q^{*}(s,a))\\
&  = Cov_{t-1}( Q^{*}(s_{t-1}, a_{t-1}), Q^{*}(s,a)) - \beta\sum_{a} \sum_{s}
\pi_{t-1}^{greedy}(a |s_{t}) F(s_{t}| s_{t-1}) Cov_{t-1}(Q^{*}(s_{t},a_{t}),
Q^{*}(s,a))\\
\end{align*}


The key is that at this step, when we compute the signal-to-noise ratio, we DO
NOT impose the Q-value equation \emph{again}, because this would imply%

\[
Q^{*}(s_{t-1}, a_{t-1}) - \beta\sum_{a} \sum_{s} \pi_{t-1}^{greedy}(a |s_{t})
f(s_{t}| s_{t-1}) Q^{*}(s_{t},a_{t}) = r(s_{s-1}, a_{t-1}),
\]
which is a known constant in the information set. By not imposing knowledge of
that equation again, and computing the signal-to-noise ratio by simply
evaluating the above covariances under the $t-1$ information set as
recursively defined is like doing only a one step update based on the Qvalue
equation, but not an infinite update. And this makes sense -- infinitely
iterating forward on the Bellman equation is solving it. We want to impose a
cost on solving it, so any deliberate iteration and attempts to solve the
Bellman equation only come through the costly signals $\eta_{t}^{R}$. \bigskip

\section{Joint actions and reasoning choices}

Conditional on a realized $s_{t}$, the objective function of the agent is to
choose an action policy $\pi(a|s_{t})$ and reasoning choices over signals
$\eta_{t}$, resulting in posterior variance of beliefs $\tilde\Sigma_{t}$, to maximize%

\[
\max_{\pi(a|s_{t}), \Sigma_{t}} \sum_{a}\widehat{Q}_{t}(a,s_{t})\pi
(a|s_{t})-w\sum_{a}\widetilde\sigma_{t}^{2}(a,s_{t})- \frac{\kappa}{2}
\ln\left[  \frac{| \Sigma_{t-1}(a,a^{\prime};s_{t})|}{|\widetilde\Sigma
_{t}(a,a^{\prime};s_{t} )|}\right]
\]
subject to two constraints
\begin{align*}
\sum_{a}\pi(a|s_{t})  &  =1\\
\underbrace{-\sum_{a}\ln(\pi(a|s_{t}))\pi(a|s_{t})}_{\equiv H_{\pi}(s_{t})}
&  \geq h\sum_{a}\widetilde\sigma_{t}^{2}(a,s_{t})
\end{align*}
where

\begin{itemize}


\item[--] $\Sigma_{t-1}(a,a^{\prime}; s_{t}) = Var(Q^{*}(a,s_{t}) |
\mathcal{I}_{t-1})$ is the beginning of period $t$ variance that agents
perceive over the vector $Q^{*}(a,s_{t})$ at a given $s_{t}$ (it is a vector
due to the vector of possible actions $a$), which variance is conditional on
the end of period $t-1$ information set $\mathcal{I}_{t-1}$. 

\item[--] $\widetilde\Sigma_{t}(a,a^{\prime};s_{t}) = \Sigma_{t-1}(a,
a^{\prime}; s_{t} | \eta_{t}) = Var(Q^{*}(a,s_{t})| \mathcal{I}_{t-1},
\eta_{t})$ is the posterior variance conditional on the time $t-1$ information
set $\mathcal{I}_{t-1}$ and also the new signals $\eta_{t}$. Note, this is not
the end-of-period $t$ posterior variance $\Sigma_{t}(a,a^{\prime}; s_{t})$
because at the end of the period you also receive and updated based on the
experience signal (discussed below). 

\item[--] $\widetilde\sigma_{t}^{2}(a,s)$ is the diagonal entry of
$\Sigma_{t-1}(a,a^{\prime};s_{t} | \eta_{t})$ corresponding to the position of
action $a$. 

\item[--] $H_{\pi}(s_{t}) \equiv-\sum_{a}\ln(\pi(a|s_{t}))\pi(a|s_{t})$ is
entropy of the actions  distribution, conditional on a state $s_{t}.$
\end{itemize}

This objective function captures:

\begin{itemize}
\item benefit of exploitation (greediness): incentive to choose the action
with highest expected value, through the linear term $\sum_{a}\widehat{Q}%
_{t}(a,s)\pi(a|s)$

\item possible disutility cost of uncertainty over the values associated to
each action (the term $w \sum_{a} \widetilde\sigma_{t}^{2}(a,s_{t}))$.
Essentially, a direct cost of cognitive dissonance, when $w>0.$ This component
going through $w$ allows us to have a benefit of reasoning that is separate
from the exploitation-experimentation tradeoff below.

\item constant marginal cost of cognitive reasoning, when $\kappa>0.$ Note:

\begin{itemize}
\item we do not model reasoning as a fixed capacity, but as a continuous choice

\item The term $\frac{1}{2} \ln\left[  \frac{|\Sigma_{t-1}(a,a^{\prime}%
;s_{t})|}{|\widetilde\Sigma_{t}(a,a^{\prime};s_{t} )|}\right] $ in the
objective function is the information content of the signals $\eta_{t}$,
measured in entropy -- i.e. that is the reduction in entropy after observing
the vector of signals $\eta_{t}$.
\end{itemize}

\item The first constraint essentially constraints $\pi(a|s_{t})$ to be proper
probability distribution (it will be associated with a Lagrange multiplier
$\lambda>0)$

\item The second constraint is an entropy regularization and its aim is to
capture the desire for exploration/experimentation. That is, as long as the
value function $Q^{*}(s,a)$ is uncertainty, there is a desire to experiment
and explore the function in other parts of the state space (exploration is
valuable here because of the model free signal agents get each period, for
free). We model this desire to experiment along the idea of ``entropy
regularization'' from ML literature, which essentially requires the entropy of
the action distribution to be bigger than some minimum threshold, thus
ensuring randomization.

Nevertheless, we recognize that there is only reason to experiment to the
extent to which the value function is uncertainty -- if $Q^{*}(a,s_{t})$ was
known, then there is no point to experimenting as there is nothing to learn.
Hence, we make the lower bound on the entropy of the action distribution
$H_{\pi}(s_{t})$ proportional to the remaining uncertainty over $Q^{*}%
(a,s_{t})$, measured by $\sum_{a}\sigma_{t}^{2}(a,s_{t})$.
\end{itemize}

\bigskip

With the LMs, objective is
\begin{align*}
&  \max_{\pi_{t}(a|s_{t}), \widetilde\Sigma_{t}}\sum_{a}\widehat{Q}%
_{t}(a,s_{t})\pi(a|s_{t})-w\sum_{a}\widetilde\sigma_{t}^{2}(a,s_{t})-
\frac{\kappa}{2} \ln\left[  \frac{|\Sigma_{t-1}(a,a^{\prime};s_{t}%
)|}{|\widetilde\Sigma_{t}(a,a^{\prime};s_{t} )|}\right] \\
&  +\delta_{t}\left[  -h\sum_{a}\widetilde\sigma_{t}^{2}(a,s_{t})-\sum_{a}%
\ln(\pi(a|s_{t}))\pi(a|s_{t})\right]  +\lambda_{t}(\sum_{a}\pi(a|s)-1)
\end{align*}


\begin{itemize}
\item \textbf{Optimal actions}. The optimal policy $\pi(a|s_{t})$ looks like a
softmax, with a state-dependent temperature%
\begin{align*}
\widehat{\pi}(a|s_{t}) =
\begin{cases}
\frac{\exp(\frac{\widehat{Q}_{t}(a,s_{t})}{\delta_{t}} )}{\sum_{a}\exp
(\frac{\widehat{Q}_{t}(a,s_{t})}{\delta_{t}})} & \text{ if } {\delta}_{t}>0\\
\widehat{\pi}^{greedy}(a|s_{t}) & \text{ if } {\delta}_{t}=0
\end{cases}
\end{align*}


where $\widehat\pi^{greedy}(a|s_{t})$ is the degenerate policy with
probability 1 on $\arg\max Q(s_{t},a)$.

To finding ${\delta}_{t}>0,$ plug in this expression into the constraint and
solve given the optimal info choice (which we characterize below)
\[
\sum_{a}\ln(\frac{\exp(\frac{\widehat{Q}_{t}(a,s)}{\widehat{\delta}_{t}}%
)}{\sum_{a}\exp(\frac{\widehat{Q}_{t}(a,s)}{\widehat{\delta}_{t}})})\frac
{\exp(\frac{\widehat{Q}_{t}(a,s)}{\widehat{\delta}_{t}})}{\sum_{a}\exp
(\frac{\widehat{Q}_{t}(a,s)}{\widehat{\delta}_{t}})}=h\sum_{a}\widetilde
{\sigma}_{t}^{2}(a,s)
\]

\end{itemize}

\textbf{Limit under }$\kappa=0.$ To get some intuition, it is useful to
consider the limit of costless reasoning $\kappa= 0.$ Formally, the objective
becomes {\small
\[
\max_{\pi_{t}(a|s_{t}), \Sigma_{t} }\sum_{a}\widehat{Q}_{t}(a,s_{t}%
)\pi(a|s_{t})-w\sum_{a}\widetilde\sigma_{t}^{2}(a,s_{t})+\delta_{t}\left[
-h\sum_{a}\widetilde\sigma_{t}^{2}(a,s_{t})-\sum_{a}\ln(\pi(a|s_{t}%
))\pi(a|s_{t})\right]  +\lambda_{t}(\sum_{a}\pi(a|s_{t})-1)
\]
}

Given costless reasoning, we will have $\widetilde\sigma_{t}^{2}(a,s_{t})=0.$
In this case, the entropy regularization constraint will not bind. So
$\delta_{t}=0$ and the optimal action follows the usual greedy policy. Further
note that $\widetilde\sigma_{t}^{2}(a,s_{t})=0$ means the agent knows the
$Q(a,s_{t})$ associated to each action, conditional on current state $s_{t}$
and conditional on following the optimal policy from now on. Thus in this
limit agent knows $Q_{\pi^{\ast}}(a, s)$ and we can replace%
\[
\widehat{Q}_{t}(a,s)=Q_{\pi^{\ast}}(a,s )
\]
Objective then is%
\[
\max_{\pi_{t}(a|s_{t})}\sum_{a}Q_{\pi^{\ast}}(a, s_{t})\pi(a|s_{t}%
)+\lambda_{t} (\sum_{a}\pi(a|s_{t})-1)
\]
which is exactly the problem solved under full rationality in the first
place:
\begin{align*}
\max_{a_{t}}Q_{\pi^{\ast}}(a_{t}, s_{t})  &  =\max_{a_{t}}\left[  u(a_{t},
s_{t})+\beta E_{t}\max_{a^{\prime}}Q_{\pi^{\ast}}(a^{\prime}, s_{t+1})\right]
\\
\pi^{\ast}(a|s)  &  =\pi^{greedy}(a|s)
\end{align*}


So, in the limit where the agent learns fully, we are back to the full
information problem as expected.

\begin{itemize}


\item \textbf{Reasoning choice.} The agent learns about the full-information
$Q_{\pi^{\ast}}(a, s).$  The agent has can generate reasoning signals
$\eta_{t}$ that are unrestricted linear functions of the underlying
$Q_{\pi^{*}}(a,s_{t})$
\[
\eta_{t} = V_{t}^{\prime}Q_{\pi^{*}}(a,s_{t}) + \varepsilon_{\eta,t}
\]
where $\varepsilon_{\eta,t} \sim N( 0, \Sigma_{\eta,t})$. The agents freely
decide on the structure of $V_{t}$ and the noise variance matrix $\Sigma
_{\eta,t}$, with only imposing the normalization that the columns of $V_{t}$
have a norm of 1.

Using the following two properties of eigenvalues,%

\[
tr(\Sigma_{t}(a,a^{\prime}; s_{t})) = \sum_{i} \lambda_{t}^{(i)}
\]
%
\[
|\Sigma_{t}(a,a^{\prime}; s_{t})| = \prod_{i} \lambda_{t}^{(i)}
\]
where $\lambda_{t}^{(i)}$ are the eigenvalues of $\widetilde\Sigma
_{t}(a,a^{\prime}; s_{t}),$ we can rewrite the objective function as
\begin{align*}
&  \max_{\pi_{t}(a|s_{t}), \Sigma_{t}}\sum_{a}\widehat{Q}_{t} (a,s_{t}%
)\pi(a|s_{t})-w\sum_{i} \lambda_{t}^{(i)} - \frac{\kappa}{2}\ln\left[
\ln(|\Sigma_{t-1}(a,a^{\prime};s_{t})|) - \sum_{i} \ln( \lambda_{t}^{(i)})
\right] \\
&  +\delta_{t}\left[  -h\sum_{i} \lambda_{t}^{(i)} -\sum_{a}\ln(\pi
(a|s_{t}))\pi(a|s_{t})\right]  +\lambda_{t}(\sum_{a}\pi(a|s)-1)
\end{align*}


The FOCs are%

\[
\lambda_{t}^{(i)} = \frac{\kappa}{w + \delta_{t} h}%
\]


Thus, we can see the interaction between actions and reasoning -- there is an
incentive to reason more when the LM $\delta_{t}$ is larger, when you would
value relaxing the  constraint on greedines induced by experimentation. The
agent can relax that  constraint by reducing uncertainy. Intuitively, the
agent values the reduction in uncertainty because that allows it to exploit
$\widehat Q_{t}(a,s_{t})$ more, in the sense of selecting an action closer to
its argmax.

Also, the above manipulations show that the solution to $\delta_{t}$ should
satisfy
\[
\sum_{a}\ln(\frac{\exp(\frac{\widehat{Q}_{t}(a,s)}{\widehat{\delta}_{t}}
)}{\sum_{a}\exp(\frac{\widehat{Q}_{t}(a,s)}{\widehat{\delta}_{t}})})\frac
{\exp(\frac{\widehat{Q}_{t}(a,s)}{\widehat{\delta}_{t}})}{\sum_{a}\exp
(\frac{\widehat{Q}_{t}(a,s)}{\widehat{\delta}_{t}})}= h \sum_{i} \lambda
_{t}^{(i)}
\]


Without loss of generality, we can assume that $\lambda_{t}^{(i)}$ are sorted
in descending order so%

\[
\lambda_{t}^{(1)} \geq\lambda_{t}^{(2)} \geq... \geq\lambda_{t}^{(N)}
\]


Importantly, by standard properties of variance matrices and conditional
variances, since $\Sigma_{\eta}$ is positive definite by assumption, the
(sorted) eigenvalues of $\Sigma_{t}(a,a^{\prime}; s_{t})$ must be weakly
smaller than those of the prior variance $\Sigma_{t-1}(a,a^{\prime};s_{t})$.
Essentially, this is the ``no forgetting'' constraint from our univariate
case, and it implies that%

\[
\lambda_{t}^{(i)} \leq\lambda_{t-1}^{(i)}
\]


So imposing the no-forgetting constraint we see that the optimal information
choice over the structure and variances of the signals $\eta_{t}$ must be such that%

\[
\lambda_{t}^{(i)} = \min\{ \frac{\kappa}{w + \delta_{t} h}, \lambda
_{t-1}^{(i)} \}
\]


Hence, agents want to only reduce eigenvalues that are already bigger than
$\kappa$. But any information flow that reduces any of the eigenvalues that
already below $\kappa$ is inefficient. Hence, the optimal signal structure is
for $V_{t}$ to be the matrix of eigenvectors of the prior variance
$\Sigma_{t-1}(a,a^{\prime}; s_{t})$, and $\Sigma_{\eta,t}$ to be a diagonal
matrix with diagonal entries $\sigma_{i,t}^{2}$ where%

\[
\sigma_{i,t}^{2} =
\begin{cases}
\frac{\kappa\lambda_{t-1}^{(i)}}{ \lambda_{t-1}^{(i)} - \kappa}, & \text{if }
\lambda_{t}^{(i)} > \kappa\\
\infty, & \text{if } \lambda_{t}^{(i)} \leq\kappa
\end{cases}
\]


Typically, the eigenvalues are declining in value quickly, so only a few of
these signals will have non-infinite noise variance. Moreover the update form
$\lambda_{t-1}^{(i)}$ to $\lambda_{t}^{(i)}$ is straightforward because it is
an update over orthogonal normal variables with a vector of orthogonal signals.
\end{itemize}

\paragraph{ End of period $t$ -- Model Free learning.}

The above joint solution produces optimal reasoning signals $\eta_{t}$, and
after observing those signals and following the optimal policy function
$\widehat\pi(a|s_{t})$, also a realized action $a_{t}$. Given that action, at
the end of the period the agent observes realized utility $u(s_{t}, a_{t})$.
This knowledge implies an ``experience'' or model-free learning signal
{\small
\begin{align*}
\eta_{t}^{Q}  &  =u(s_{t},a_{t})+\beta\sum_{a}\sum_{s_{t+1}}\pi_{t}%
^{greedy}(a|s_{t+1})\widehat{Q}_{t}(s_{t+1},a)F(s_{t+1}|s_{t})\\
&  =u(s_{t},a_{t})+\beta\sum_{a}\sum_{s_{t+1}}\left[  \pi_{t} ^{greedy}%
(a|s_{t+1})\widehat{Q}_{t}(s_{t+1},a)+\pi^{\ast}(a|s_{t+1})Q_{\pi^{\ast}
}(s_{t+1},a)-\pi^{\ast}(a|s_{t+1})Q_{\pi^{\ast}}(s_{t+1},a)\right]  F(s_{t+1}
|s_{t})\\
&  =Q^{\ast}(s_{t},a_{t})+\beta\sum_{a}\sum_{s_{t+1}}\left[  \pi_{t}%
^{greedy}(a|s_{t+1})\widehat{Q}_{t}(s_{t+1},a)-\pi^{\ast}(a|s_{t+1}
)Q_{\pi^{\ast}}(s_{t+1},a)\right]  F(s_{t+1}|s_{t})
\end{align*}
} where $\pi_{t}^{greedy}(a|s_{t+1}) = \arg\max_{a} \widehat Q_{t}(s_{t+1}%
,a)$. Because the greedy policies are degenerate, the last line simplifies to%
\[
\eta_{t}^{Q}=Q^{\ast}(s_{t},a_{t})+\beta\sum_{s_{t+1}}\left[  \widehat{Q}%
_{t}(s_{t+1}, \pi_{t}^{greedy}(a|s_{t+1}) )- Q_{\pi^{\ast}}(s_{t+1},\pi^{\ast
}(a|s_{t+1}))\right]  F(s_{t+1}|s_{t})
\]


The agent has beliefs about $Q_{\pi^{\ast}}(s_{t},a)$ for all actions $a$ but
does not know $\pi^{\ast}(a|s_{t+1})$. To start with, let us approximate by
setting $\pi^{\ast}(a|s_{t+1})=\pi_{t}^{greedy}(a|s_{t+1})$
\[
\eta_{t}^{Q}=Q^{\ast}(s_{t},a_{t})+\beta\sum_{s_{t+1}}\left[  \widehat{Q}%
_{t}(s_{t+1}, \pi_{t}^{greedy}(a|s_{t+1}) )- Q_{\pi^{\ast}}(s_{t+1}, \pi
_{t}^{greedy}(a|s_{t+1}) )\right]  F(s_{t+1}|s_{t})
\]
and use this signal to update beliefs about $Q(s,a)$ as per formulas above.
Note: the approximation $\pi^{\ast}(a|s_{t})=\pi_{t-1}^{greedy}(a|s_{t})$
introduces another layer of noise, which we could potentially add in the
future, but for now let's ignore it.

Using this signal, we update the beliefs of the agent one last time to obtain
the end of period $t$ beliefs, after both active deliberation and experiential learning.

\section{Application: firms' investment}

Consider a firm $j^{\prime}$s investment problem. Starts the period with some
capital $k_{j,t}.$ Budget constraint says dividends equal
\[
d_{j,t}=z_{j,t}k_{j,t}^{\alpha}+i_{j,t}\left[  1-\phi(i_{j,t},k_{j,t})\right]
\]
where $z_{j,t}k_{j,t}^{\alpha}$ is output of firm, $z_{j,t}$ is a technology
shock,  $i_{j,t}$ is investment and $\phi(i_{j,t},k_{j,t})$ could be some
investment adjustment costs. For example, it might be the usual $i/k$ form:%
\[
\phi(i_{j,t},k_{j,t})=\varkappa\frac{i_{j,t}}{k_{j,t}}%
\]
Capital accumulation:
\[
k_{j,t+1}=(1-\delta)k_{j,t}+i_{j,t}%
\]
and a law of motion for technology $F(z_{j,t+1}|z_{j,t}),$ for example an
AR(1). 

Firm's problem is
\[
V(k_{j,t},z_{j,t})=\max_{i_{j,t}}\left[  u(d_{j,t})+\beta E_{t}V(k_{j,t+1}%
,z_{j,t+1})\right]
\]
where $u(d_{j,t})$ is current payoff function, over dividends. We can have
linear or concave $u(.),$ for example $u(d)=\log(d).$

For now, we make entry and exit completely exogenous. With some probability
$\gamma,$ an individual firm $j$ exits the economy. In particular, a new firm
replaces it. The new firm will have some time-zero prior beliefs over the
objects defined in previous sections (and further below), and have some
starting capital $k_{j,0}=k_{0}.$ We will endogenize at some point more the
entry and exit decisions. 

Because all firms will solve separately their problem there is no need to keep
track of the $j$ index for what follows. 

Call the firm's state $s_{t}=\left(  k_{t},z_{t}\right)  .$ We would then need
to: 

1. solve for the fully rational policy function $i^{\ast}(s_{t}).$ 

2. Implement our learning model, as described above. Recall the optimal
Q-function, where the action is now investment
\[
Q_{i^{\ast}}(s_{t},i_{t})=r(s_{t},i_{t})+\beta E_{t}Q_{i^{\ast}}%
(s_{t+1},i^{\ast}(s_{t+1}))
\]
where $r(s_{t},i_{t})=u\left(  z_{t}k_{t}^{\alpha}+i_{t}-\phi(i_{t}%
,k_{t})\right)  .$ The agent will learn about the optimal Q-value function
$Q_{\pi^{\ast}}(s,i)$ defined above.

\bigskip

To start with we could even make it simpler, but shutting down stochastic
shock $z_{t}.$ Instead it's entirely about transitional dynamics, through the
endogenous state $k_{t}.$
\end{document}